\newpage
\section{Kesimpulan dan Saran}

\subsubsection{Hasil Decision Tree tanpa hyperparameter tuning}
Berikut adalah hasil confusion matrix dari model Decision Tree yang telah diuji pada data uji, Gambar \ref{fig:decision_tree} menunjukkan confusion matrix yang diperoleh dari model Decision Tree tanpa melakukan hyperparameter tuning. Sehingga baseline model Decision Tree ini akan menjadi acuan untuk perbandingan dengan model-model lainnya yang akan diuji pada tahap selanjutnya. Namun kita tetap perlu menambahkan beberapa setting manual karena pada bab sebelumnya kita menemukan imbalance data pada dataset yang digunakan. Oleh karena itu, kita perlu mengatur beberapa parameter pada model Decision Tree untuk mengatasi masalah tersebut. Berikut adalah parameter yang digunakan pada model Decision Tree:
